% Part: proof-theory
% Chapter: sequent-calculus
% Section: interpretation-rules

\documentclass[../../../include/open-logic-section]{subfiles}

\begin{document}

\olfileid{pt}{seq}{irl}

\olsection{বিধিগুলির ব্যাখ্যা}

সিকোয়েন্টের ব্যাখ্যা মনে করি:
$\Gamma \Sequent \Delta$
!!a{structure}~$\Struct{M}$-তে
সত্য যদি হয় পূর্বপক্ষের
অন্তত একটি $!A \in \Gamma$ মিথ্যা,
নয়তো উত্তরপক্ষের
অন্তত একটি~$!A \in \Delta$ সত্য।
\Log{G3c}-র যৌক্তিক বিধিগুলি
তাদের প্রধান !!{formula}s-এর
সত্যশর্ত সংকেতায়িত করে।
\RightR{\lif} দেখি। এর প্রধান
!!{formula} $!A \lif !B$ সত্য
যদি $!A$ মিথ্যা অথবা~$!B$
সত্য হয়। তাই $\ \Sequent !A \lif !B$
তখনই সত্য, যখন $!A \Sequent !B$
সত্য। এই সিকোয়েন্টগুলির
পূর্বপক্ষ ও উত্তরপক্ষে প্রসঙ্গের
!!{formula}s~$\Gamma$, $\Delta$
যোগ করলে ঠিক \RightR{\lif}-এর
সিদ্ধান্ত ও পূর্বধারণা পাই।
\LeftR{\lif}-এর ক্ষেত্রেও
অনুরূপ। $!A \lif !B$ মিথ্যা
যদি $!A$~সত্য ও $!B$~মিথ্যা
হয়। তাই $!A \lif !B
\Sequent \ $ তখনই সত্য, যখন
$\ \Sequent !A$ এবং $!B \Sequent \ $
উভয়ই সত্য। \LeftR{\lif}-এর
সিদ্ধান্ত তার পূর্বধারণাগুলির
যৌগের সমতুল্য। এই ধরন
\Log{G3c}-র বিধিগুলির বিশুদ্ধতা
নিশ্চিত করে: সব পূর্বধারণা
সত্য হলে সিদ্ধান্তও সত্য।
এগুলি পূর্ণতাও নিশ্চিত করে।

আসলে যেকোনো সত্যমান-অপেক্ষকধর্মী
সংযোজকের সত্যশর্ত এভাবে কয়েকটি
সিকোয়েন্টে বলা যায়। কারণ
চিরায়ত যুক্তিবিদ্যার প্রতিটি
সত্যমান-অপেক্ষককে যৌগীয়
স্বাভাবিক রূপের !!a{formula}
দিয়ে প্রকাশ করা যায়: আণবিক
ও নঞর্থিত আণবিক !!{formula}s-এর
বিযোজনগুলির যৌগ। যেমন,
$!A \oplus !B$ তখনই সত্য ধরি,
যখন $!A$ ও $!B$-এর ঠিক একটি
সত্য—“বর্জনমূলক অথবা”।
তাহলে $!A \oplus !B$ তখনই
সত্য, যখন $(!A \lor !B) \land
(\lnot !A \lor \lnot !B)$ সত্য;
আর $(\lnot !A \lor !B)
\land (!A \lor \lnot !B)$
সত্য হলে সেটি মিথ্যা। তাই
$\oplus$-এর সিকোয়েন্ট-কলনের
বিধি দেওয়া যায় এভাবে:
\begin{align*}
&
\Axiom$\Gamma \fCenter \Delta, !A, !B$
\Axiom$!A, !B, \Gamma \fCenter \Delta$
\RightLabel{\RightR{\oplus}}
\BinaryInf$\Gamma \fCenter \Delta, !A \oplus !B$
\DisplayProof
\\
& \Axiom$!A, \Gamma \fCenter \Delta, !B$
\Axiom$!B, \Gamma \fCenter \Delta, !A$
\RightLabel{\LeftR{\oplus}}
\BinaryInf$!A \oplus !B, \Gamma \fCenter \Delta$
\DisplayProof
\end{align*}
এই বিধিগুলিও বিশুদ্ধ ও পূর্ণ
হবে। প্রতিটির সিদ্ধান্ত তখনই
সত্য, যখন তার সব পূর্বধারণা সত্য।
% BN-SRC-627 TARGET-CORRECTION-BEGIN
\par\smallskip\noindent\textnormal{\small\textbf{উৎসসংশোধন (BN-SRC-627):} XOR-এর দ্বিতীয় বিধির সিদ্ধান্তে \( !A\oplus !B \) বাম দিকে, কিন্তু উৎসে লেবেল \(\RightR{\oplus}\) পুনরাবৃত্ত। বিধিটির লেবেল \(\LeftR{\oplus}\) করা হয়েছে; দুই পূর্বধারণা ও সিদ্ধান্ত অপরিবর্তিত।}
% BN-SRC-627 TARGET-CORRECTION-END

\begin{prob}
ধরি, $!A \downarrow !B$ তখনই
সত্য, যখন $!A$ বা~$!B$ কোনোটিই
সত্য নয় (NOR)। \Log{G3c}-তে
এর বিধিগুলি কী হবে? (এমন দুটি
সিকোয়েন্ট বের করো, যারা একসঙ্গে
তখনই সত্য, যখন $!A \downarrow !B$
সত্য। তাতে \RightR{\downarrow}
বিধি পাবে। তারপর এমন একটি
সিকোয়েন্ট বের করো, যা তখনই
সত্য, যখন $!A \downarrow !B$
মিথ্যা। তাতে \LeftR{\downarrow}
বিধি পাবে।)
\end{prob}

\Log{G3c}-তে প্রত্যেক
সংযোজক ও পরিমাণসূচকের ঠিক দুটি
বিধি: একটি বাঁ, একটি ডান।
কিন্তু \Log{G1c}-তে
\LeftR{\land}-এর দুটি এবং
\RightR{\lor}-এর দুটি বিধি।
কারণ \Log{G1c}-র আদর্শ
জেন্টজেনের মূল ব্যবস্থা~\Log{LK}-র
একটি প্রকরণ~\Log{LJ} ছিল
গুরুত্বপূর্ণ অচিরায়ত স্বজ্ঞাবাদী
যুক্তির জন্য। স্বজ্ঞাবাদী
যুক্তির বিশুদ্ধ ও পূর্ণ ব্যবস্থা
পেতে \Log{LJ}-র প্রতিটি
সিকোয়েন্টের উত্তরপক্ষে সর্বাধিক
একটি !!{formula} রাখা হয়।
তাই \Log{G3c}-র
\RightR{\lor} বিধি ব্যবহার
করা যায় না: তার পূর্বধারণার
উত্তরপক্ষে $!A$ ও~$!B$
উভয়ই আছে। জেন্টজেন তাই
\Log{LJ}-তে \RightR{\lor}-এর
দুটি বিধি ব্যবহার করেন এবং
\Log{LK}-তেও সেগুলি রাখেন।
একইভাবে স্বজ্ঞাবাদী
প্রকরণ~\Log{G1i} হলো
\Log{G1c}, যেখানে প্রত্যেক
উত্তরপক্ষে সর্বাধিক একটি
!!{formula} রাখা হয়।
(\Log{G3c}-রও স্বজ্ঞাবাদী
প্রকরণ আছে, তবে তার কিছু
বিধি \Log{G3c}-র সংশ্লিষ্ট
বিধির থেকে আলাদা। যেমন,
সেখানেও \RightR{\lor}-এর
দুটি বিধি।)

\Log{G1c}-র গঠনগত
বিধিগুলির বিশুদ্ধতা সহজে
বোঝা যায়। যেমন,
$\Gamma \Sequent \Delta$-র
বাঁ দিকে মিথ্যা !!{formula}
অথবা ডান দিকে সত্য !!{formula}
থাকলে $\Gamma \Sequent
\Delta, !A$-তেও সেই
একই সাক্ষ্য থাকে; $!A$
সত্য না মিথ্যা, তাতে
কিছু যায় আসে না।
তাই \RightR{\Weakening}-এর
পূর্বধারণা সত্য হলে সিদ্ধান্তও
সত্য। এখানে বিপরীতটি সত্য
নয়: $!A$ সত্য হওয়ার
কারণেই সিদ্ধান্ত সত্য হতে
পারে, অথচ পূর্বধারণা সত্য
নাও হতে পারে।
% BN-SRC-628 TARGET-CORRECTION-BEGIN
\par\smallskip\noindent\textnormal{\small\textbf{উৎসসংশোধন (BN-SRC-628):} সিকোয়েন্টের শুরুর সত্যশর্ত অনুযায়ী বাঁ দিকে মিথ্যা অথবা ডান দিকে সত্য সূত্র থাকলে সিকোয়েন্ট সত্য। উৎসের গঠনগত-বিধি অনুচ্ছেদে এই দুই পাশ উলটো লেখা ছিল; সত্যতার সাক্ষ্য সঠিক পাশে রাখা হয়েছে।}
% BN-SRC-628 TARGET-CORRECTION-END

তাহলে গঠনগত বিধিগুলি আদৌ কেন দরকার? উদাহরণস্বরূপ,
$\ \Sequent !A \lor \lnot !A$ !!{prove} করতে
সংকোচন (\RightR{\Contraction}, \LeftR{\Contraction})
দরকার। $!A \Sequent !A$ থেকে শুরু করলে
প্রথমে \RightR{\lnot} প্রয়োগ করে
$\Sequent !A, \lnot !A$ পাওয়া যায়। তারপর
\RightR{\lor}-এর দুই রূপ পরপর প্রয়োগ করা যায়:
একবার $\lnot !A$ থেকে $!A \lor \lnot !A$
পেতে, আর একবার $!A$ থেকে $!A \lor \lnot !A$
পেতে। এতে পাওয়া যায়
$\ \Sequent !A \lor \lnot !A, !A \lor \lnot !A$।
অতএব \Log{G1c}-তে $\Sequent !A \lor \lnot !A$
পেতে হলে !!a{proof}-এ একই !!{formula}-র
দুটি উপস্থিতিকে একটি উপস্থিতিতে ``সংকুচিত''
করতে পারতে হবে:
\[
\Axiom$!A \fCenter !A$
\RightLabel{\RightR{\lnot}}
\UnaryInf$\fCenter !A, \lnot !A$
\RightLabel{\RightR{\lor}}
\UnaryInf$\fCenter !A, !A \lor \lnot !A$
\RightLabel{\RightR{\lor}}
\UnaryInf$\fCenter !A \lor \lnot !A, !A \lor \lnot !A$
\RightLabel{\RightR{\Contraction}}
\UnaryInf$\fCenter !A \lor \lnot !A$
\DisplayProof
\]
বস্তুত এমন বৈধ সিকোয়েন্ট আছে, যেগুলি
দুর্বলীকরণ বা সংকোচন ছাড়া \Log{G1c}-তে
!!{prove} করা যায় না। যেমন,
$!A \Sequent !B \lif !A$ সিকোয়েন্টের
!!a{proof} \Log{G1c}-তে
\RightR{\lif} দিয়ে শেষ হতে হবে
(কারণ কেবল $!B \lif !A$-ই
মুখ্য হতে পারে)। এর জন্য
$!B, !A \Sequent !A$ পূর্বধারণা
দরকার। সেটি আবার $!A \Sequent !A$
স্বতঃসিদ্ধ থেকে পেতে
\LeftR{\Weakening} প্রয়োজন:
\[
\Axiom$!A \fCenter !A$
\RightLabel{\LeftR{\Weakening}}
\UnaryInf$!B, !A \fCenter !A$
\RightLabel{\RightR{\lif}}
\UnaryInf$!A \fCenter !B \lif !A$
\DisplayProof
\]

\Log{G3c}-তে সংকোচন ও দুর্বলীকরণের
কোনো বিধিই দরকার হয় না। দুর্বলীকরণ
স্বতঃসিদ্ধগুলির মধ্যেই অন্তর্নিহিত।
দুই পূর্বধারণাবিশিষ্ট বিধিতে
$\Gamma$, $\Delta$ প্রসঙ্গের
দুটি অনুলিপি একটিতে মিলিয়ে
সংকোচন অধিকাংশ ক্ষেত্রে এড়ানো যায়।
\Log{G1c} ও \Log{G3c} উভয়
পদ্ধতিতেই এ ধরনের ``একই প্রসঙ্গ
ভাগ করে নেওয়া'' বিধি আছে।
এক পূর্বধারণাবিশিষ্ট বিধিগুলির মধ্যে
সংকোচন কেবল \RightR{\lor},
\LeftR{\land}, \RightR{\lexists}
ও~\LeftR{\lforall}-এর জন্য দরকার।
\Log{G3c}-তে এই বিধিগুলির রূপ
সংকোচনকে সম্পূর্ণ অপ্রয়োজনীয় করে।
\Log{G2c} নামে আরেকটি পদ্ধতি
আছে (\olref{tab:G2c} দেখুন),
যার দুই পূর্বধারণাবিশিষ্ট বিধিগুলিতে
প্রসঙ্গ দুটি স্বাধীন। \Log{G2c}-তে
সংকোচন \Log{G1c}-র চেয়েও বেশি
গুরুত্বপূর্ণ।

\subfile{rules-G2c}

\end{document}
